\section{Methods}

\subsection{Decision models and queries}
Two decision models were used, Jev~1.13 (TypeSafe System One~\cite{jev}) and d3-ultra (vLLM Semantic Router~\cite{d3ultra}), both through their web services and with the same requests. A request describes the case in short text fields and asks one or more questions, each with a fixed set of allowed answers. The model returns a probability for each allowed answer, or a score between 0 and 1. Neither service offers settings such as a temperature or a random seed, and repeated calls gave stable answers: asked twice, 200 drug-claim pairs received the same verdict in 99.5\% of cases from Jev and in all cases from d3-ultra. The tests were first run with Jev, and d3-ultra was added afterwards with the same questions. All answers were stored, so every result can be rebuilt without calling the models again. The exact questions, a worked example for each test and the code are in the repository.

\subsection{Study screening}
The data were 39 systematic reviews from SYNERGY+~\cite{synergy}, each with titles, abstracts, eligibility criteria and final inclusion decisions. The records of each review were split at random into parts for development and for testing, and all results come from the test parts (18,183 records, 577 of them eligible). For each record the models read the title, the abstract and the review's eligibility criteria and scored whether the record should go on to full-text review. This score ordered the records. The pre-set cut-off of a review was the highest score that still found 95\% of the eligible studies among its development records, and it was applied unchanged to the test records.

The comparator was ASReview LAB~3.0.8~\cite{asreview}, run once per review in simulation mode from one eligible and one ineligible record. The measure is work saved at 95\% recall~\cite{cohen2006}: the share of records left unread at the point where 95\% of the eligible records have been found, minus 0.05.

\subsection{Drug-trial claims}
DrugCon-Bench~\cite{lan2026} pairs claims from randomised-trial abstracts on the same drug and disease. It holds 1,731 pairs, of which 542 contradict each other. The benchmark's labels started from pairs flagged by one of the published models, which annotators then checked and extended. Each pair was put to the models as one question with two answers, that the claims contradict or that they can coexist, worded after the benchmark authors' own prompt. A pair was called contradictory when that answer had the higher probability, a rule fixed before any run. The measure is F$_1$ for the contradictory pairs, computed with the authors' evaluation code. The scores of the eleven published models were taken from the benchmark paper, using for each model the better of its two prompts. They were published without intervals and are treated as fixed values.

\subsection{Cell types}
Expert-annotated cells of pancreas and heart came from Tabula Sapiens~\cite{tabula}, and a fifth of each cell type was held out for testing (2,440 pancreas and 7,075 heart cells). The candidate genes were the CellMarker~2.0~\cite{cellmarker} markers of each tissue's cell types. For each gene the models were asked which of the tissue's cell types it marks, or none of them, and the returned probabilities form a table of genes by cell types that takes the place of the database list. Each cell then received the type whose marker genes carried the largest share of its expression, by the same rule for the model and database lists.

Applied in this way, the database list assigned almost all heart cells to one type, so for heart it was also scored in two further ways. The z-score version standardises each gene's expression across the test cells before scoring. The size-normalised version, added after the first results were seen, divides each type's score by the number of genes listed for it. Each model built its list three times; the text reports the average, and comparisons with the database use Jev's first build. Figure~\ref{fig:cell} also shows a second arm in which genes often mentioned together with each gene in the literature were added to the question.

\subsection{Mass spectrometry}
One run of a human cell line from a public dataset~\cite{chick} was searched with Sage~\cite{sage}, using standard tryptic settings, against the human proteins together with the proteins of a plant, \emph{Arabidopsis thaliana}, which cannot be in the sample. The search gave 37,061 matches to real sequences. Each match was put to the models as one yes-or-no question listing the search engine's numbers for that match, such as scores, mass errors, matched fragments and explained signal. The models never saw the peptide sequence, the species or any re-scorer output, and the probability of ``yes'' ranked the matches.

The comparator was mokapot~0.10.0~\cite{mokapot}, trained on the run with its standard inputs. Two further checks used the same run: a classifier trained on only the numbers the models saw, and mokapot given each model's probability as one more input. For the small-run test, mokapot was retrained on random subsets of 500 to 20,000 matches, 20 per size.

The error among accepted matches was estimated by entrapment~\cite{wen}. A match to a plant protein must be wrong, so the number of such matches, scaled for the sizes of the two protein sets, estimates how many accepted human matches are wrong. Each method's matches were accepted down to the score at which this estimate reached 1\%.

All comparisons are paired. Intervals are 95\% bootstrap intervals, obtained by resampling the unit that carries the dependence: the review for screening, the group of pairs that share abstracts for drug claims, the donor for heart (approximate with only six donors~\cite{field2007}) and the match for mass spectrometry. The pancreas data carry no donor labels, so that interval was computed over cells. A difference is called not detectable when its interval includes zero, which does not show that the methods are equal. Intervals are descriptive and were not adjusted for the number of comparisons. The stored results behind every comparison, and the code that computes its interval, are in the repository.
All comparisons are paired. Intervals are 95\% bootstrap intervals, obtained by resampling the unit that carries the dependence: the review for screening, the group of pairs that share abstracts for drug claims, the donor for heart (approximate with only six donors~\cite{field2007}) and the match for mass spectrometry. The pancreas data carry no donor labels, so that interval was computed over cells. A difference is called not detectable when its interval includes zero, which does not show that the methods are equal. Intervals are descriptive and were not adjusted for the number of comparisons. The values and intervals of every comparison are in the repository.
